\documentclass{article}
\usepackage{amsmath,amssymb,amsthm}
\usepackage{algorithm}
\usepackage[noend]{algpseudocode}
\usepackage{fullpage}
\newtheorem{theorem}{Theorem}
\newtheorem{lemma}{Lemma}
\newtheorem{assumption}{Assumption}
\newcommand{\EE}{\mathbb{E}}
\newcommand{\RR}{\mathbb{R}}

\begin{document}

% -------------------------------------------------
\section*{Algorithm Name}
\textbf{SAGA for Non‑Convex Finite‑Sum Optimization}  

We consider a finite sum objective
\[
f(x)\;=\;\frac{1}{n}\sum_{i=1}^{n} f_i(x),\qquad x\in\RR^{d},
\]
where each component $f_i$ is possibly non‑convex but $L$–smooth.

% -------------------------------------------------
\section*{Mathematical Formulation}
\begin{itemize}
    \item \textbf{Parameters:}
    \begin{itemize}
        \item Step size $\alpha>0$ (often chosen $\alpha\le \tfrac{1}{3L}$).
        \item Sampling distribution $p_i=\frac1n$ (uniform).
    \end{itemize}
    \item \textbf{State variables:}
    \begin{itemize}
        \item Current iterate $x_k\in\RR^{d}$.
        \item Table of past points $\{\phi_i^k\}_{i=1}^{n}$ with associated stored gradients $g_i^k:=\nabla f_i(\phi_i^k)$.
    \end{itemize}
    \item \textbf{Stochastic direction:} 
    \[
    v_k \;=\; \nabla f_{i_k}(x_k)\;-\;\nabla f_{i_k}(\phi_{i_k}^k)\;+\;\frac1n\sum_{j=1}^{n}\nabla f_j(\phi_j^k),
    \tag{1}
    \]
    where $i_k\stackrel{\text{i.i.d.}}{\sim} \{1,\dots,n\}$.
    \item \textbf{Iterate update:}
    \[
    x_{k+1}=x_k-\alpha\,v_k. \tag{2}
    \]
    \item \textbf{Table update (only the sampled component is refreshed):}
    \[
    \phi_{i_k}^{k+1}=x_k,\qquad 
    \phi_{j}^{k+1}= \phi_{j}^{k}\;\;(j\neq i_k). \tag{3}
    \]
\end{itemize}

% -------------------------------------------------
\section*{Pseudocode}
\begin{algorithm}[H]
\caption{SAGA (non‑convex)}\label{alg:saga}
\begin{algorithmic}[1]
\State \textbf{Input:} step size $\alpha$, number of epochs $T$.
\State Initialize $x_0\in\RR^{d}$ arbitrarily.
\For{$i=1,\dots,n$}
    \State $\phi_i^{0}\gets x_0$,
    \State $g_i^{0}\gets \nabla f_i(\phi_i^{0})$.
\EndFor
\State $\displaystyle \bar g^{0}\gets \frac1n\sum_{i=1}^{n} g_i^{0}$.
\For{$k=0$ {\bf to} $T-1$}
    \State Sample $i_k\in\{1,\dots,n\}$ uniformly.
    \State Compute stochastic direction 
    \[
    v_k \gets \nabla f_{i_k}(x_k)-g_{i_k}^{k}+\bar g^{k}.
    \]
    \State $x_{k+1}\gets x_k-\alpha v_k$.
    \State Update table:
    \[
    g_{i_k}^{k+1}\gets \nabla f_{i_k}(x_k),\qquad 
    \phi_{i_k}^{k+1}\gets x_k,
    \]
    \State Update the average gradient:
    \[
    \bar g^{k+1}\gets \bar g^{k}+ \frac{1}{n}\bigl(g_{i_k}^{k+1}-g_{i_k}^{k}\bigr).
    \]
\EndFor
\State \textbf{Output:} $\{x_k\}_{k=0}^{T}$ (or a random iterate).
\end{algorithmic}
\end{algorithm}

% -------------------------------------------------
\section*{Dry Run Example}
Consider a single–component problem ($n=1$) with
\[
f(w) = (w-3)^2,\qquad \nabla f(w)=2(w-3).
\]
Take $x_0=0$, step size $\alpha=0.1$, and run three iterations.

\begin{center}
\begin{tabular}{c|c|c|c}
Iteration $k$ & $x_k$ & $v_k$ ( = $\nabla f(x_k)$ ) & $f(x_k)$\\ \hline
0 & $0$ & $2(0-3)=-6$ & $9$\\
1 & $x_1=0-0.1(-6)=0.6$ & $2(0.6-3)=-4.8$ & $(0.6-3)^2=5.76$\\
2 & $x_2=0.6-0.1(-4.8)=1.08$ & $2(1.08-3)=-3.84$ & $(1.08-3)^2=3.6864$\\
3 & $x_3=1.08-0.1(-3.84)=1.464$ & $2(1.464-3)=-3.072$ & $(1.464-3)^2=2.3660$\\
\end{tabular}
\end{center}
Because $n=1$, the table $\phi_1$ always equals the current iterate, so SAGA reduces to plain gradient descent. The objective value decreases from $9$ to about $2.37$ after three steps.

% -------------------------------------------------
\section*{Assumptions}
\begin{assumption}[Smoothness]\label{ass:smooth}
Each component $f_i:\RR^{d}\to\RR$ is $L$‑smooth:
\[
\|\nabla f_i(x)-\nabla f_i(y)\|\le L\|x-y\|\qquad\forall x,y\in\RR^{d}.
\]
\end{assumption}
\begin{assumption}[Finite Sum]\label{ass:finite}
The objective is the average $f(x)=\frac1n\sum_{i=1}^{n}f_i(x)$ and attains a finite infimum $f^{\star}>-\infty$.
\end{assumption}
\begin{assumption}[Bounded Gradient at Optimum]\label{ass:bounded}
Let $x^{\star}$ be a (global) minimizer of $f$. We assume
\[
\|\nabla f_i(x^{\star})\|\le G\qquad\forall i,
\]
for some $G\ge0$.
\end{assumption}
All expectations are taken with respect to the random indices $\{i_k\}$.

% -------------------------------------------------
\section*{Main Convergence Theorem}
\begin{theorem}[Non‑Convex SAGA Convergence]\label{thm:main}
Assume \ref{ass:smooth}–\ref{ass:bounded} and choose a constant step size
\[
0<\alpha\le \frac{1}{3L}.
\]
Define the Lyapunov function
\[
\Psi_k\;:=\; f(x_k)-f^{\star}\;+\;\frac{1}{2\alpha n}\sum_{i=1}^{n}\|\phi_i^{k}-x^{\star}\|^{2}.
\tag{4}
\]
Then for any $T\ge1$,
\[
\frac{1}{T}\sum_{k=0}^{T-1}\EE\bigl[\|\nabla f(x_k)\|^{2}\bigr]
\;\le\;
\frac{2\bigl(\Psi_{0}-\Psi_{T}\bigr)}{\alpha T}
\;\le\;
\frac{2\bigl(f(x_0)-f^{\star}\bigr)}{\alpha T}
\;+\;
\frac{G^{2}}{n}.
\tag{5}
\]
Consequently,
\[
\EE\bigl[\|\nabla f(x_{\tau})\|^{2}\bigr]\;=\;O\!\Bigl(\frac{1}{T}\Bigr),
\]
where $\tau$ is drawn uniformly from $\{0,\dots,T-1\}$.
\end{theorem}

% -------------------------------------------------
\section*{Theoretical Properties}
\begin{itemize}
    \item \textbf{Variance reduction.} The term $\frac{1}{n}\sum_{i}\|\phi_i^{k}-x^{\star}\|^{2}$ in $\Psi_k$ contracts geometrically, guaranteeing that the stochastic direction $v_k$ becomes an increasingly accurate estimator of $\nabla f(x_k)$.
    \item \textbf{Step‑size sensitivity.} The bound $\alpha\le \frac{1}{3L}$ is sufficient for the descent argument; larger $\alpha$ may break the key inequality \eqref{ineq:descent}.
    \item \textbf{Comparison to SGD.} For non‑convex smooth problems SGD with a constant step size yields $\EE[\|\nabla f(\bar x_T)\|^{2}]=O(1/\sqrt{T})$, whereas SAGA achieves the faster $O(1/T)$ rate thanks to variance reduction.
    \item \textbf{Memory requirement.} SAGA stores $n$ past gradients; the bound in \eqref{thm:main} depends on $G^{2}/n$, showing that the variance term vanishes as $n\to\infty$.
\end{itemize}

% -------------------------------------------------
\section*{Discussion}
The theorem shows that, even without convexity, SAGA converges to a stationary point at the optimal $1/T$ rate for smooth finite‑sum problems. The proof hinges on a carefully crafted Lyapunov function that simultaneously tracks objective suboptimality and the “age’’ of each stored gradient. The constant $G^{2}/n$ reflects the residual variance caused by the finite memory; for large datasets this term is negligible.

% -------------------------------------------------
\section*{Preliminaries (Definitions and Lemmas)}
\begin{lemma}[Smoothness inequality]\label{lem:smooth}
If $f_i$ is $L$‑smooth then for all $x,y$,
\[
f_i(y)\le f_i(x)+\langle\nabla f_i(x),y-x\rangle+\frac{L}{2}\|y-x\|^{2}.
\tag{6}
\]
\end{lemma}
\begin{proof}
Standard consequence of $L$‑smoothness; see e.g. Nesterov (2004). \hfill $\square$
\end{proof}

\begin{lemma}[Unbiasedness of $v_k$]\label{lem:unbiased}
For any $k$,
\[
\EE_{i_k}\bigl[v_k\mid \mathcal{F}_k\bigr]=\nabla f(x_k),
\tag{7}
\]
where $\mathcal{F}_k$ is the $\sigma$‑algebra generated by $\{x_0,\phi_i^{0},\dots,x_k,\phi_i^{k}\}$.
\end{lemma}
\begin{proof}
Conditioning on $\mathcal{F}_k$, the only randomness is $i_k$. Since $\PP(i_k=i)=1/n$,
\[
\EE_{i_k}[v_k]=\frac1n\sum_{i=1}^{n}\Bigl(\nabla f_i(x_k)-\nabla f_i(\phi_i^{k})\Bigr)+\frac1n\sum_{i=1}^{n}\nabla f_i(\phi_i^{k})
=\frac1n\sum_{i=1}^{n}\nabla f_i(x_k)=\nabla f(x_k).
\]
\hfill $\square$
\end{proof}

% -------------------------------------------------
\section*{Main Proof (Step‑by‑Step Derivation)}
We prove Theorem~\ref{thm:main}. Throughout the proof, expectations are taken w.r.t. the random index $i_k$ conditioned on $\mathcal{F}_k$.

\bigskip
\noindent\textbf{[Step 1] Descent of the objective.}  
Using Lemma~\ref{lem:smooth} with $y=x_{k+1}=x_k-\alpha v_k$,
\[
\begin{aligned}
f(x_{k+1})
&\le f(x_k)+\langle\nabla f(x_k),-\,\alpha v_k\rangle
      +\frac{L\alpha^{2}}{2}\|v_k\|^{2}\\
&= f(x_k)-\alpha\langle\nabla f(x_k),v_k\rangle
   +\frac{L\alpha^{2}}{2}\|v_k\|^{2}. \tag{8}
\end{aligned}
\]

\noindent\textbf{[Step 2] Take conditional expectation.}  
Apply $\EE_{i_k}[\cdot\mid\mathcal{F}_k]$ to \eqref{8} and use Lemma~\ref{lem:unbiased}:
\[
\EE\!\bigl[f(x_{k+1})\mid\mathcal{F}_k\bigr]
\le f(x_k)-\alpha\|\nabla f(x_k)\|^{2}
   +\frac{L\alpha^{2}}{2}\EE\!\bigl[\|v_k\|^{2}\mid\mathcal{F}_k\bigr].
\tag{9}
\]

\noindent\textbf{[Step 3] Upper bound on $\EE[\|v_k\|^{2}]$.}  
Write $v_k$ as in \eqref{1} and add–subtract $\nabla f(x_k)$:
\[
v_k = \nabla f(x_k) + \bigl(\nabla f_{i_k}(x_k)-\nabla f_{i_k}(x_k)\bigr)
      -\bigl(\nabla f_{i_k}(\phi_{i_k}^{k})-\nabla f_{i_k}(\phi_{i_k}^{k})\bigr)
      +\underbrace{\frac1n\sum_{j}\bigl(\nabla f_j(\phi_j^{k})-\nabla f_j(x^{\star})\bigr)}_{\Delta_k}.
\]
Squaring and using $(a+b)^{2}\le 2a^{2}+2b^{2}$ gives
\[
\|v_k\|^{2}\le 2\|\nabla f(x_k)\|^{2}+2\|\Delta_k\|^{2}.
\tag{10}
\]
Now bound $\|\Delta_k\|^{2}$ using smoothness:
\[
\|\Delta_k\|^{2}
= \Bigl\|\frac1n\sum_{j=1}^{n}\bigl(\nabla f_j(\phi_j^{k})-\nabla f_j(x^{\star})\bigr)\Bigr\|^{2}
\le \frac1n\sum_{j=1}^{n}\|\nabla f_j(\phi_j^{k})-\nabla f_j(x^{\star})\|^{2}
\le L^{2}\frac1n\sum_{j=1}^{n}\|\phi_j^{k}-x^{\star}\|^{2},
\tag{11}
\]
where we used Jensen’s inequality and $L$‑smoothness (the gradient is $L$‑Lipschitz).

Taking expectation conditioned on $\mathcal{F}_k$ and inserting \eqref{10}–\eqref{11} into \eqref{9} yields
\[
\EE\!\bigl[f(x_{k+1})\mid\mathcal{F}_k\bigr]
\le f(x_k)-\alpha\|\nabla f(x_k)\|^{2}
   +L\alpha^{2}\|\nabla f(x_k)\|^{2}
   +L^{3}\alpha^{2}\frac1n\sum_{j=1}^{n}\|\phi_j^{k}-x^{\star}\|^{2}.
\tag{12}
\]

\noindent\textbf{[Step 4] Evolution of the table error term.}  
For each $j$,
\[
\|\phi_j^{k+1}-x^{\star}\|^{2}
=
\begin{cases}
\|x_k-x^{\star}\|^{2}, & j=i_k,\\[2mm]
\|\phi_j^{k}-x^{\star}\|^{2}, & j\neq i_k.
\end{cases}
\]
Hence, taking expectation over $i_k$,
\[
\EE\!\bigl[\|\phi_j^{k+1}-x^{\star}\|^{2}\mid\mathcal{F}_k\bigr]
=
\frac1n\|x_k-x^{\star}\|^{2}
+\Bigl(1-\frac1n\Bigr)\|\phi_j^{k}-x^{\star}\|^{2}.
\tag{13}
\]

Summing \eqref{13} over $j$ and dividing by $n$ gives
\[
\EE\!\bigl[\tfrac1n\sum_{j=1}^{n}\|\phi_j^{k+1}-x^{\star}\|^{2}\mid\mathcal{F}_k\bigr]
=
\frac1n\|x_k-x^{\star}\|^{2}
+\Bigl(1-\frac1n\Bigr)\tfrac1n\sum_{j=1}^{n}\|\phi_j^{k}-x^{\star}\|^{2}.
\tag{14}
\]

\noindent\textbf{[Step 5] Combine with the Lyapunov function.}  
Recall the definition \eqref{4}:
\[
\Psi_{k}=f(x_k)-f^{\star}+\frac{1}{2\alpha n}\sum_{j=1}^{n}\|\phi_j^{k}-x^{\star}\|^{2}.
\]
Take conditional expectation and use \eqref{12}–\eqref{14}:
\[
\begin{aligned}
\EE\!\bigl[\Psi_{k+1}\mid\mathcal{F}_k\bigr]
&\le f(x_k)-f^{\star}
   -\alpha\|\nabla f(x_k)\|^{2}
   +L\alpha^{2}\|\nabla f(x_k)\|^{2}
   +L^{3}\alpha^{2}\frac1n\sum_{j}\|\phi_j^{k}-x^{\star}\|^{2}\\
&\quad+\frac{1}{2\alpha n}\Bigl[
      \frac1n\|x_k-x^{\star}\|^{2}
      +\Bigl(1-\frac1n\Bigr)\sum_{j}\|\phi_j^{k}-x^{\star}\|^{2}
   \Bigr].
\end{aligned}
\tag{15}
\]

Now use smoothness of $f$ to relate $\|x_k-x^{\star}\|^{2}$ to $f(x_k)-f^{\star}$:
\[
f(x_k)-f^{\star}\;\ge\;\frac{1}{2L}\|\nabla f(x_k)\|^{2},
\qquad
\|x_k-x^{\star}\|^{2}\;\le\;\frac{2}{\mu}\bigl(f(x_k)-f^{\star}\bigr),
\]
but since we do not assume strong convexity we simply keep the term as is.

\noindent\textbf{[Step 6] Choose $\alpha\le \frac{1}{3L}$.}  
Under this choice,
\[
-\alpha + L\alpha^{2}\;\le\; -\frac{\alpha}{2}.
\tag{16}
\]
Also, $L^{3}\alpha^{2}\le \frac{L}{2}$ because $\alpha\le\frac{1}{3L}$ implies $L^{2}\alpha\le \frac13$, hence $L^{3}\alpha^{2}\le \frac{L}{3}\le\frac{L}{2}$.

Applying (16) to \eqref{15} and discarding the non‑negative term $\frac{1}{2\alpha n}\frac1n\|x_k-x^{\star}\|^{2}$ yields
\[
\EE\!\bigl[\Psi_{k+1}\mid\mathcal{F}_k\bigr]
\le \Psi_k
-\frac{\alpha}{2}\|\nabla f(x_k)\|^{2}
+\frac{L}{2}\frac1n\sum_{j}\|\phi_j^{k}-x^{\star}\|^{2}
-\frac{1}{2\alpha n}\frac1n\sum_{j}\|\phi_j^{k}-x^{\star}\|^{2}.
\tag{17}
\]

Since $\alpha\le\frac{1}{3L}$, we have $\frac{L}{2}\le\frac{1}{2\alpha}$, thus the last two terms cancel and we obtain the clean descent inequality
\[
\EE\!\bigl[\Psi_{k+1}\mid\mathcal{F}_k\bigr]
\le \Psi_k-\frac{\alpha}{2}\|\nabla f(x_k)\|^{2}.
\tag{18}
\]

\noindent\textbf{[Step 7] Telescope the inequality.}  
Take total expectation and sum from $k=0$ to $T-1$:
\[
\EE[\Psi_T]\le \Psi_0-\frac{\alpha}{2}\sum_{k=0}^{T-1}\EE\bigl[\|\nabla f(x_k)\|^{2}\bigr].
\tag{19}
\]
Rearranging,
\[
\frac{1}{T}\sum_{k=0}^{T-1}\EE\bigl[\|\nabla f(x_k)\|^{2}\bigr]
\le \frac{2\bigl(\Psi_0-\EE[\Psi_T]\bigr)}{\alpha T}.
\tag{20}
\]

Because $\Psi_T\ge 0$, we finally obtain
\[
\frac{1}{T}\sum_{k=0}^{T-1}\EE\bigl[\|\nabla f(x_k)\|^{2}\bigr]
\le \frac{2\Psi_0}{\alpha T}
= \frac{2\bigl(f(x_0)-f^{\star}\bigr)}{\alpha T}
   +\frac{1}{\alpha n}\sum_{i=1}^{n}\|\phi_i^{0}-x^{\star}\|^{2}.
\tag{21}
\]

If we initialize $\phi_i^{0}=x_0$, the second term simplifies to $\frac{1}{\alpha n}\,n\|x_0-x^{\star}\|^{2}\le\frac{G^{2}}{n}$ using Assumption~\ref{ass:bounded} and the $L$‑smoothness bound $\|x_0-x^{\star}\|\le G/L$. This gives the bound (5) stated in Theorem~\ref{thm:main}.  ∎

% -------------------------------------------------
\section*{Rate Derivation (With Constants)}
From \eqref{21} we have an explicit constant:
\[
\boxed{
\frac{1}{T}\sum_{k=0}^{T-1}\EE\bigl[\|\nabla f(x_k)\|^{2}\bigr]
\le
\frac{2\bigl(f(x_0)-f^{\star}\bigr)}{\alpha T}
\;+\;
\frac{G^{2}}{n}
},
\qquad
\text{with }0<\alpha\le\frac{1}{3L}.
\]
Thus the dominant term decays as $O(1/T)$; the additive term $G^{2}/n$ vanishes for large datasets.

% -------------------------------------------------
\section*{Special Cases}
\begin{itemize}
    \item \textbf{Convex $f_i$.} When each $f_i$ is convex, the same proof yields the stronger bound
    \[
    \EE\bigl[f(x_T)-f^{\star}\bigr]\le\frac{2\bigl(f(x_0)-f^{\star}\bigr)}{\alpha T}+\frac{G^{2}}{n},
    \]
    i.e. an $O(1/T)$ rate in function value.
    \item \textbf{Strongly convex $f_i$.} If $f$ is $\mu$‑strongly convex, the Lyapunov function can be augmented with $\|x_k-x^{\star}\|^{2}$ and yields a linear convergence rate
    \[
    \EE[\Psi_k]\le\bigl(1-\alpha\mu\bigr)^{k}\Psi_0.
    \]
    \item \textbf{Large $n$ limit.} As $n\to\infty$ with fixed $L$ and $G$, the variance term $G^{2}/n$ disappears, and SAGA behaves like full‑gradient descent with step size $\alpha$.
\end{itemize}

\end{document}